\documentclass[11pt]{article}
% =====================================================================================================================
% main.tex -- Universities, not departments: where academic prestige meets graduate pay (Haotian Yang).
% Main text in Nature Human Behaviour format. Figures: outputs/figures/wapman2026/ (scripts/83_nexus_figures.py and
% scripts/86_nexus_figures_extra.py). SI pointers resolve through xr against si.aux: compile si.tex first, then
% main.tex (latexmk -pdf si; latexmk -pdf main).
% =====================================================================================================================
\usepackage[a4paper,margin=2.5cm]{geometry}
\usepackage[T1]{fontenc}
\usepackage{times}
\usepackage[scaled=0.92]{helvet}
\usepackage{amsmath,amssymb}
\usepackage{graphicx}
\usepackage{booktabs}
\usepackage{array}
\usepackage{xcolor}
\usepackage{caption}
\DeclareCaptionLabelSeparator{bar}{\space\textbar\space}
\captionsetup{font=small,labelfont=bf,labelsep=bar}
\usepackage[super,sort&compress,comma]{natbib}
\usepackage{lineno}
\usepackage{xr-hyper}
\usepackage{hyperref}
\hypersetup{colorlinks=true,linkcolor=blue!50!black,citecolor=blue!50!black,urlcolor=blue!50!black}
\externaldocument[SI-]{si}

\linespread{1.213}
\renewcommand{\topfraction}{0.9}
\renewcommand{\bottomfraction}{0.8}
\renewcommand{\textfraction}{0.07}
\renewcommand{\floatpagefraction}{0.8}
\setcounter{topnumber}{3}
\setcounter{totalnumber}{4}
\setlength{\emergencystretch}{3em}
\renewcommand{\figurename}{Fig.}

\newcommand{\ci}[2]{$[#1,\allowbreak\,#2]$}
\newcommand{\SIsec}[1]{Supplementary Information, section~\ref{SI-#1}}
\newcommand{\SItab}[1]{Supplementary Information, Table~\ref{SI-#1}}
\newcommand{\figfile}[1]{../../outputs/figures/wapman2026/#1}


\title{\bfseries Universities, not departments: where academic prestige meets graduate pay}
\author{Haotian Yang\textsuperscript{1,*}\\[6pt]
\small \textsuperscript{1}Bocconi University, Milan, Italy\\
\small \textsuperscript{*}Corresponding author: Haotian Yang, \href{mailto:haotian.yang@studbocconi.it}{haotian.yang@studbocconi.it}}
\date{}

\begin{document}
\maketitle
\linenumbers

\begin{abstract}
\noindent Subject rankings rate a university's programme in a field partly by its academic standing, yet it is unclear
whether that standing goes with what graduates earn, and whether through the department or the whole university. We
compare the two orderings for the same university $\times$ field programmes in the United States, using departments'
positions in the faculty hiring network, and in the United Kingdom, using the national research assessment. Academic
standing and graduates' pay agree substantially across universities but much less among the departments of one
university: the correlation is 0.45 against 0.05 in the US and 0.62 against 0.10 in the UK, a contrast that survives
allowing for noisier within-university measures. Within US universities, doctoral departments' standing still tracks
research output. We argue that academia and the labour market share mainly a university-level ordering, so a
department's standing says little about its graduates' pay once its university is known.
\end{abstract}

\medskip
% =====================================================================================================================
% Introduction (no heading, as Nature Human Behaviour requires)

Universities are ranked as wholes and by subject. A subject ranking rates a university's programme in one field,
economics at one university or nursing at another, and it does so partly by surveying academics and employers
\citep{qs_subject}. Rankings shift where students apply \citep{luca2013salience}, so what they reward matters to the
applicants who read them. Academic standing can also be read from behaviour rather than from surveys: departments
reveal their judgement of one another in which doctoral programmes they hire faculty from, and these hires trace a
steep prestige hierarchy of departments \citep{clauset2015systematic, wapman2022quantifying}. Outside academia, the
outcome that applicants and rankings most often invoke is what a programme's graduates earn.

Most evidence on how graduates' pay depends on where they studied, however, compares universities as wholes or fields
as wholes. Among students admitted to similar colleges, the pay advantage of the more selective one is small on
average \citep{dale2002estimating}. Returns differ by field, and attending a more selective institution pays
relatively little by comparison \citep{kirkeboen2016field}. In the UK, no subject's league-table rating correlates
positively and significantly with returns once selectivity is held fixed \citep{britton2022how}. None of this work
asks whether academic standing goes with pay among the departments inside one university. This, we argue, leaves
open the question that decides what a subject ranking can tell an applicant: does the academic standing that
corresponds to graduates' pay belong to the university or to the department? There is reason to expect the answer to
depend on the level. Hiring committees rank departments field by field, whereas employers see the university's name
on every graduate's record and can pool their experience with all of its graduates, much as buyers pool experience
across the products of one brand \citep{wernerfelt1988umbrella, tadelis1999name}.

In this work, we compare two orderings of the same university $\times$ field programmes, their academic standing and
their graduates' pay, in the United States and the United Kingdom. In the US, a department's academic standing is its
rank in its field's faculty-hiring network; the UK has no such rank, so there it is the department's rating in the
national research assessment. We measure how closely the two orderings agree across the universities that teach a
field and among the departments of a single university. Because the second comparison is the noisier one, we test
the difference against simulated data that reproduce the noise of each measure. We then ask whether the department's
rank still carries academic information inside its university, how much pay goes with it there, and whether it
improves out-of-sample predictions of pay. We also follow agreement over graduates' careers and in UK subjects paid on
national scales.

We find that academic standing and pay agree substantially across universities but far less within them: their
correlation is 0.45 across universities and 0.05 within them in the US, and 0.62 against 0.10 in the UK. In
simulations that reproduce each measure's noise, equal agreement at both levels never yields a within-university
correlation as low as observed. Inside a university, the department's rank is not noise, since it tracks research
output far more closely than pay, and in an exploratory decomposition what it adds to predictions of pay is almost
all information about its university. The academic hierarchy and the labour market, we argue, share mainly a ranking
of universities. For applicants, and for rankings that score departments partly by academic reputation, a
department's academic standing therefore says little about its graduates' pay once its university is known.

% =====================================================================================================================
\section*{Results}

\subsection*{Two orderings of the same programmes}

A programme is a university's bachelor's degree in one field, such as economics at one university, and we order the
same programmes in two ways. The first ordering is academic standing. In the US we use two published ranks from the
faculty-hiring network \citep{wapman2022quantifying}: $F$, a department's rank among the departments of its field,
and $G$, its university's rank in the network that pools all fields. The UK has no department-level hiring rank, so
there the department's standing is its rating in the Research Excellence Framework (REF) 2021
(ref.~\citenum{ref2021results}). The second ordering is graduates' median pay: four years after completion in the US College
Scorecard \citep{scorecard}, with the Census Post-Secondary Employment Outcomes (PSEO) following fixed cohorts for ten
years \citep{pseo2026}, and one to five years after graduation in UK tax records \citep{dfe_leo_provider}.

We refer to the correlation of the two orderings as their \emph{agreement} and compute it in three ways on the same
programmes. Agreement across universities compares the programmes of one field at different universities, so it
mixes differences between universities with differences between the departments of one university. Agreement
between universities correlates university averages. Agreement within universities compares the departments of one
university with each other, after university and field fixed effects. Put differently, a university whose departments
all rank high and whose graduates all earn well raises agreement across and between universities but leaves
agreement within them unchanged. The three are different statistics, so the gap between them is not a share of
agreement explained by the university (Methods).

If pay is set competitively and employers learn about graduates symmetrically, as our model assumes
(\SIsec{sec:si-model}), pay reflects employers' beliefs about a programme
\citep{farber1996learning, altonji2001employer, macleod2017big}, and agreement can be read as how far two markets,
employers and academic hiring committees, rank programmes alike \citep{macleod2015reputation}. We use this reading to
interpret the results rather than to predict them.

\subsection*{Universities and departments}

Fig.~\ref{fig:levels} sets agreement across universities beside agreement within them, both computed on normal
scores so that the two orderings share one scale in every field (Methods). We find that the two levels
differ sharply. In the US, $F$ and pay agree at 0.45 across universities but at only 0.05
within them, although both statistics use the same programmes and only the fixed effects separate them. The UK, where
the academic measure is itself departmental, gives the same picture: REF rating and pay agree at 0.62 across
universities and at 0.10 within them. Because the across statistic mixes variation between universities with
variation within them, it understates agreement about universities proper: agreement between universities, on
university means, is 0.69 in the US and 0.78 in the UK (\SIsec{sec:si-contrast}). Intake attainment, a third ordering
that shares the university-wide components of the other two \citep{macleod2017big}, shows the same contrast with both
REF rating and pay.

\begin{figure}[!htbp]
\centering
\includegraphics[width=8.7cm]{\figfile{nexus_fig1}}
\caption{\textbf{Academic standing and pay agree across universities and far less within them.} Points are observed
correlations with 95\% intervals: blue compares programmes across universities, vermilion the departments of one
university. Grey bands show the within-university correlation simulated with equal agreement at both levels. Intake
rows are supporting evidence. US and UK programmes; details in \SIsec{sec:si-figs}.}\label{fig:levels}
\end{figure}

Two qualifications concern what the within statistic isolates. It holds the field fixed, so it says nothing about the
value of choosing one field over another, or about comparing universities within a field. Nor is it purely
departmental: where a university trait counts for more in some fields than in others, part of that trait survives the
university fixed effects. Netting out field-specific slopes on selectivity and on $G$ removes part of it and brings
the US within value close to zero (\SIsec{sec:si-contrast}).

A lower correlation within universities is to be expected even if agreement were equal at both levels, however,
because differences between the departments of one university are measured with more noise. Within UK universities,
for example, REF ratings from 2014 and 2021 correlate at only about 0.5 (\SItab{tab:si-reliab}). We therefore ask how
low the within correlation would be if academic standing and pay agreed equally at both levels. To this end, we
simulate data sets on the real layout of programmes, universities and fields, with equal latent agreement between and
within universities and with each measure's noise set to reproduce its observed reliability, and we pass every
simulated data set through the same pipeline as the data (Methods). Under equal agreement, the US within
correlation would be about 0.31 (\SItab{tab:si-calib}). None of the simulated values is as low as the observed 0.05,
and the same holds for every UK pair (Fig.~\ref{fig:levels}, grey bands). To check the test itself, we re-ran the
whole procedure, nuisance estimation included, on data simulated under equal agreement. It rejects at about the
nominal rate when the noise model is right, and up to about five times too often when noise sits at the pessimistic
end of its range. Even that distortion is small beside the observed gap. The contrast also holds on the largest US sample, which does not need selectivity data,
and with REF 2014 as a re-measurement of UK academic standing (\SIsec{sec:si-audit}). When the simulation is instead fitted with
separate agreement at each level, it puts the latent agreement of $F$ and pay at about 0.08 within universities against
0.70 between them (\SItab{tab:si-model}). These model-based values rest on the latent model, whereas the observed
correlations do not, and we use them only to illustrate magnitudes.

The contrast for intake is less robust than the one for REF rating and pay. When we rebuild the UK analysis on
programmes with at least one high-attainment entrant, or on programmes with many graduates in the earnings data, the
reliability rule written into the script before its estimate (Methods) still separates the levels for REF rating and
pay but no longer for some pairs involving intake (\SItab{tab:si-robust}). Because that rule proved miscalibrated,
these subsets describe rather than test. The like-for-like within correlations are all positive, so the contrast is
one of degree. We read it as follows: academic standing and pay do agree about departments, but most of what they
share is an ordering of universities.

\subsection*{Academic information within universities}

Is the department's rank simply noise once its university is fixed? We find that it is not, at least for doctoral
departments. Within universities, departments ranked higher in hiring publish more per faculty member in the National
Research Council's assessment of doctoral programmes \citep{nrc2011data}, and this association is far stronger than
their association with graduates' pay on the same programmes, which is close to zero (Fig.~\ref{fig:academic}A); it
holds net of department size and doctoral production (\SIsec{sec:si-pcp}). For these departments, the weak within-university agreement with pay
therefore does not arise because the rank is noise, although noise in the pay medians could still hide a relation with
pay. Nor is the hiring rank mainly a count of doctorates: a rank that allows for each
department's volume of doctoral production and hiring keeps all of the field-level agreement with pay
(\SIsec{sec:si-slope}).

\begin{figure}[!htbp]
\centering
\includegraphics[width=8.7cm]{\figfile{nexus_fig_academic}}
\caption{\textbf{Inside a university, the department's rank tracks research more than pay, and its position among the
university's departments adds little to predicting pay.} (A) Within-university correlations of the hiring rank $F$ with
NRC research measures (green) and with pay on the same programmes (vermilion). (B) Gain in out-of-sample $R^2$ from
$F$ and its parts; rows below the rule are exploratory. Bars: 95\% intervals. Details in \SIsec{sec:si-figs}.}\label{fig:academic}
\end{figure}

\subsection*{What the department adds}

We first compare a department's own rank with its university's rank as predictors of pay across universities. Across
US fields, agreement of $F$ with pay is on average slightly lower than agreement of $G$ with pay, and correcting $F$
for its greater noise shrinks this gap without closing it (\SItab{tab:si-map}). $F$ thus predicts pay no better than
$G$ across universities. That does not make it redundant: within universities, where the university fixed effects
hold $G$ fixed, $F$ still goes with pay, as the positive within correlation of Fig.~\ref{fig:levels} shows.

On the scale of pay, we find the department's part to be small (Fig.~\ref{fig:gradient}). Across the universities of a
field, graduates of a department one within-field standard deviation (SD) higher in $F$ earn 0.077 log points more four
years after completion. Among the departments of one university, the same gradient is 0.010 log points, i.e., about
1\% of pay. One SD is a sizeable step: in a typical field it spans about 36 places of the hiring rank. On the
programmes with selectivity data, a specification that adds field-specific slopes on mean SAT, admission rate, Pell
share and $G$ puts the gradient at about 0.6\%, with an interval that reaches below zero; the samples differ, so
these numbers are not a decomposition. Adding the department's doctoral production to that specification barely changes it.
The controls absorb most of $F$, so what remains of $F$ is noisy, but the complete hiring records let us measure that
noise, and correcting for it raises the gradient only to about 0.8\% per SD. Since US selectivity is measured for the
university, the gradient also carries any sorting of stronger students into a university's higher-ranked departments.

\begin{figure}[!htbp]
\centering
\includegraphics[width=8.7cm]{\figfile{nexus_fig_gradient}}
\caption{\textbf{Within a university, the department's pay gradient is small.} Change in log median pay per within-field SD
of academic standing, (A) across and within US universities and (B) within universities, by specification, for the US
hiring rank and the UK REF rating. Bars: 95\% intervals (dashed: cluster-robust); markers without bars have no interval.
Values and samples in \SItab{tab:si-gradient}.}\label{fig:gradient}
\end{figure}

The UK gradient, which rests on REF ratings rather than on hiring, is of the same order, although the data cannot
rank it against the US one. Within a university, a subject whose department stands one SD higher in REF 2021 rating
has graduates who earn about 1.8\% more five years after graduation, and about 1.2\% more with subject-specific slopes on $G$ and
selectivity, the analogue of the US main specification. REF 2014 assessed research while these cohorts studied,
whereas REF 2021 assessed it largely after they had been taught, and with REF 2014 the gradient is larger, about
3.2\%. Because the UK data also record each programme's own intake, we can ask how much of this gradient moves with
who enrols: the point estimate is about 60\%, too imprecise to read as a decomposition of the gradient
(\SItab{tab:si-gradient}).

Out of sample, the department's rank does improve prediction. Adding it with field-specific slopes raises the
out-of-sample $R^2$ for the programme pay of a held-out university by about 1.4 points, more than the threshold we had
set for calling the department's contribution small. Field-specific slopes on the university's own standing,
however, give the same gain without $F$. In an exploratory decomposition, the department's position among its
university's departments adds about nothing for a new university and about a tenth of a point within a known one
(Fig.~\ref{fig:academic}B). Once its university's standing is known, then, a department's place
among its university's departments adds little to the prediction of its graduates' pay.

\subsection*{The field map}

Agreement also differs from field to field. Measured in each US field as the Spearman correlation between $F$ and
four-year median pay across the universities that teach the field, it is positive in most fields
(Fig.~\ref{fig:map}A) and averages 0.38, a field-by-field statistic that is not directly comparable with the pooled
values of Fig.~\ref{fig:levels}. Among the fields taught at many universities, it ranges from 0.04 in music to 0.70 in
computer engineering, with computer science at 0.68 and nursing at 0.17. Part of this spread is noise in the pay
medians, but not all of it: propagating each median's sampling noise leaves most of the heterogeneity and nearly the
same ordering of fields (\SIsec{sec:si-data}), and PSEO earnings order fields much as the Scorecard does
(\SItab{tab:si-map}).

\begin{figure}[!htbp]
\centering
\includegraphics[width=\textwidth]{\figfile{nexus_fig2}}
\caption{\textbf{Agreement varies by field.} (A) US: correlation across universities of the hiring rank $F$ with four-year
pay, by field; marker area grows with the number of universities; bars: $\pm1.96$ bootstrap SE.
(B) UK: the same for the university-wide rank $G$ and five-year pay; orange marks subjects paid on national scales.
Top rows: means with 95\% intervals. Details in \SIsec{sec:si-figs}.}\label{fig:map}
\end{figure}

The field map does not, however, show where agreement sits. Fields in which pay follows $F$ closely are the fields in
which pay follows university selectivity closely: across the larger US fields, these two field-level correlations have
a true-score correlation of 0.87 (\SItab{tab:si-map}). The UK offers a check with the university-wide rank $G$, built
from the academic careers that researchers report themselves, which orders UK universities much as REF does
(Methods). Agreement of $G$ with five-year pay varies substantially across subjects
(Fig.~\ref{fig:map}B), and it is high in the subjects in which pay correlates strongly with university selectivity
(\SItab{tab:si-map}). In both countries, then, a field with high agreement could have a steep gradient between
departments or simply a steep gradient between universities. The field map alone cannot tell the two apart; the
two-level contrast above can.

Agreement about universities is itself partly, but not wholly, shared with intake. In the US, controlling for mean
SAT, admission rate and other university traits lowers mean field-level agreement, on the fields with SAT data, from
0.38 to 0.10, and dating these controls to graduates' entry years leaves the net agreement about where it was
(\SItab{tab:si-map}). The UK data control graduates' own prior attainment directly: among graduates in the same
attainment band, about two thirds of the agreement of $G$ with pay survives (\SIsec{sec:si-slope}). The two controls
remove different things, because university-level selectivity also absorbs peer and resource effects
\citep{dale2002estimating}.

\subsection*{Careers}

The remaining results are secondary: they concern agreement across universities and rest on stronger
assumptions. We observe that agreement across universities grows as careers advance. Within fixed US graduation
cohorts, agreement of $F$ with median pay rises by about 0.03 per year since graduation (Fig.~\ref{fig:career}A). In
the UK, agreement of $G$ with pay rises at a similar rate, from 0.38 one year after graduation to 0.48 five years after
(Fig.~\ref{fig:career}C), and it does so also when pay is standardised for graduates' prior attainment. Part of the US
rise may come from graduate study rather than from careers: if graduates of higher-ranked universities more often
leave the earnings data to study and return with higher pay, agreement rises mechanically. Controlling for the share
of a university's bachelor's graduates who later earn a research doctorate absorbs almost two thirds of the slope.
That share, however, is closely related to $G$, so the control removes university standing as well, and the data
cannot say how much of the absorbed part is graduate study.

\begin{figure}[!htbp]
\centering
\includegraphics[width=\textwidth]{\figfile{nexus_fig3}}
\caption{\textbf{Agreement grows with years since graduation, and the growth appears in the loading on the university-wide
rank.} (A) US mean agreement of $F$ (blue) and $G$ (pink) with pay, fixed cohorts; (B) rank-regression loadings of
pay on $F$ and $G$; (C) UK mean agreement of $G$ with pay. Bars show 95\% intervals; PSEO and LEO pay. Details in
\SIsec{sec:si-figs}.}\label{fig:career}
\end{figure}

A change within a cohort also mixes career time with calendar time. If calendar and cohort trends are not negative,
the part due to career time lies between about 0.015 and 0.03 per year, i.e., roughly half of the observed rise or
more (\SIsec{sec:si-career}). Nor does the rise mainly reflect the census division in which graduates work. Net of the
wage level of those divisions, five-year agreement keeps its size, and about a sixth of the career rise moves with
destination wages, a quarter with destination wages and the out-of-state share together (\SIsec{sec:si-career}).

Which part of the academic ordering gains weight? $F$ and $G$ are highly correlated (\SItab{tab:si-map}), and their
agreement with pay rises in parallel (Fig.~\ref{fig:career}A). Rank regressions of pay on both separate them: the
loading on $G$ grows by about 0.034 per year, while the loading on $F$ shows no detectable trend
(Fig.~\ref{fig:career}B). Whether the two trends differ from each other depends on the estimator, so we do not claim a
difference between them; net of selectivity, the partial correlation of $G$ with pay still rises, although the
corresponding regression loading is not distinguishable from zero (\SIsec{sec:si-career}). We thus find no sign that
the department's own rank, net of its university's, gains weight over the first ten years, although the data cannot
rule out a rise over nine years about as large as the $F$ loading at one year (about 0.13). This statement concerns the loadings, not the within-university
correlation of Fig.~\ref{fig:levels}, which we do not track over careers. In the model, learning about individual
graduates leaves the covariance of pay with university standing unchanged at names that employers see often, so a
rising loading on $G$ with no detectable trend in the loading on $F$ points to pay growth that follows university standing or to learning
about rarely seen names (\SIsec{sec:si-model}). Besides the graduate study noted above, the timing of the ranks offers
a competing explanation: they describe the faculty of 2011--2020, closer in time to these cohorts' later earnings than
to their earliest.

\subsection*{National pay scales}

In UK subjects whose graduates are paid largely on national scales, pay is set partly by schedule rather than by
employers. Three subjects fall in this group: nursing and midwifery, medicine and dentistry, and education and
teaching. In these, agreement of $G$ with five-year pay averages 0.12, in contrast to 0.52 in the other subjects
(Fig.~\ref{fig:map}B), and the gap remains when pay is standardised for graduates' prior attainment. The pattern
parallels the finding of \citet{britton2022how} that UK returns barely follow selectivity in medicine and education,
which they attribute to centrally regulated wages. With three subjects, however, it marks a descriptive boundary rather
than an identified mechanism, and the classification by subject is coarse: allied health, whose graduates are partly
paid on national bands, agrees about as strongly as the other subjects (\SItab{tab:si-map}). A pay scale also need not
lower agreement by itself. In the model, a schedule that only rescales market pay leaves correlations unchanged
(\SIsec{sec:si-model}), so weak agreement under national scales can reflect noisier medians, schedule variation
unrelated to standing, or weaker agreement in the underlying market, and these data do not separate the three.

% =====================================================================================================================
\section*{Discussion}

Academic standing and graduates' pay order the same programmes in a common way, but mainly through the university.
The two orderings agree substantially across universities and far less among the departments of one university, in
two countries with different academic measures, and the contrast survives a test that allows for the noisier
within-university comparison. The weak agreement within universities does not mean the rank lacks academic content:
among NRC-assessed doctoral departments it tracks research output inside a university far more closely than pay. In an
exploratory decomposition, almost all of what the department's rank adds to out-of-sample predictions of pay is
information about its university, and fields differ in agreement largely as they differ in how steeply pay follows
university selectivity. In the UK, intake attainment shows the same contrast, less robustly, and over careers agreement
across universities grows, a rise detectable in the loading on the university-wide rank but not in that on the
department's rank. Read together, these results suggest that what the
academic hierarchy and the labour market share is mainly a ranking of universities, with a small departmental
component. To the question of which organisational level of academic standing corresponds to graduates'
pay, our answer is: mostly the university.

The pattern fits two findings in economics. \citet{macleod2017big} find that employers reward a college's
reputation, measured by its graduates' admission scores, and that this reward falls when individual skill becomes
visible. UK
league-table ratings of subjects show no significant positive association with returns once selectivity is held
fixed \citep{britton2022how}. Our results add a revealed academic ordering and locate where it meets pay: mainly at
the university. For the study of faculty-hiring hierarchies \citep{clauset2015systematic, wapman2022quantifying, morgan2022socioeconomic},
they indicate that the US hiring rank carries information about graduates' pay mostly through its university-wide
component, which it shares with selectivity. Studies that use US faculty-hiring ranks as proxies for programme quality
outside academia should therefore treat their department-level component with caution.

Why agreement within universities is weak is not identified, but the model narrows the possible readings
(\SIsec{sec:si-model}). If pay is set competitively and employers learn about graduates symmetrically, pay reflects
employers' beliefs, and the within-university covariance of standing and pay is proportional to the product of two
quantities: $\kappa$, the correlation of departments' academic standing with their graduates' department-specific
productivity, and $\Lambda$, which grows with how far employers attribute productivity to departments and learn it
from graduates' own records. Weak attribution and slow learning therefore shrink the covariance and the pay gradient.
In the correlation, however, they act through the reliability of pay, which we measure and build into the test, and
through a factor $q$: the reliability-corrected within correlation equals $\kappa\sqrt{q}$, where $q$ lies between
employers' attribution weight and 1. This factor is small only if employers attribute pay to departments from noisy
records while learning almost nothing within four years. With learning as fast as estimated elsewhere
\citep{lange2007speed, arcidiacono2010beyond}, $q$ is close to 1. Within the model, and if employers' prior uncertainty
about departments matches the true spread of departments, the small model-based within correlations then point to a
small $\kappa$: departments' academic standing would track the department-specific productivity of their graduates
only weakly. The positive control fits this reading, since within universities the department's rank tracks research
output far more closely than pay. Outside the model, the REF's focus on research could also lower the measured within correlation.

Two implications follow, within the limits set out below. First, rankings that compare departments partly on academic
reputation \citep{qs_subject} may, as guides to pay, largely restate the standing of the university. Second, for an
applicant comparing programmes, a department's academic rank says little about its graduates' pay relative to other
departments of the same university: one within-field SD of department standing goes with about 0.5\% to 1\% higher
pay in the US and 1\% to 3\% in the UK. These figures describe who earns what, not what a programme does to a student's
pay.

Our design is descriptive. Agreement can arise from common causes, and the public aggregates
we use do not separate selection, value added and employer screening; in particular, the selectivity controls do not
rule out sorting on traits they do not measure. We observe graduates' earnings, not employers' judgements, and reading
pay as employers' valuation requires the model's assumptions of competitive pay and symmetric learning. Medians miss
the upper tail of pay, where elite returns concentrate \citep{zimmerman2019elite, chetty2023diversifying}, and
Scorecard medians cover only federally aided completers who work; Scorecard also reports one set of medians for
campuses that share an identifier, which pools some branch campuses, although dropping them leaves the contrast in
place (\SIsec{sec:si-data}). The academic measures differ by country: the US hiring rank describes the faculty of
2011--2020, the UK department result rests on REF ratings of research rather than on hiring, and the UK
university-wide rank uses self-reported academic careers. Because $F$ and $G$ come from overlapping hiring data, their
errors may be correlated, which affects the comparisons of $F$ with $G$ and the joint career loadings, though not the
within-university contrast. The positive control covers only doctoral departments assessed by the NRC, which sit in
the upper part of the hiring-rank distribution, and it does not rule out that noise in the pay medians hides a
within-university relation with pay. Intake records who enrols, not what students prefer, so the intake results describe
sorting, not choices. Graduate-study timing and the timing of the ranks cannot be excluded as sources of the career
rise, the career-time ranges are identified sets under sign restrictions, and pay is netted of destination wages only
between census divisions.

The calibrated test is a parametric bootstrap whose nuisance parameters are fitted on the same data, and its main
simulation gives each university one effect common to all its fields; it rejects equal agreement under the fitted
measurement model. The rejection rests on that latent structure rather than on the test's calibration. If university
standing loads differently on different fields, a null chosen after a development run to reproduce the observed
field-to-field spread of agreement still rejects for the academic--pay pairs in both countries. The US rejection fails,
on one of three loading families, once that spread is pushed about a quarter beyond what we observe: it holds at 1.21
and fails at 1.29 times the observed spread. The UK academic--pay rejection fails only at two to three times the
observed spread. We did not audit heterogeneous loadings and misspecified noise jointly, and under heterogeneous
loadings the intake pairs support no conclusion. The model-based correlations rest on the same latent model with classical errors, and their
approximate intervals leave out the uncertainty of the fitted parameters. The reliability rule we wrote into the script
before its estimate proved miscalibrated, so the UK robustness subsets that rely on it are descriptive. The calibrated test,
the like-for-like US contrast and the production adjustment inside the gradient design were all added after the first
results had been seen. The calibration audit, the positive control and the prediction analysis were specified in
their scripts after the main results had been seen, and the decomposition of the prediction gain after its own
results. Tests that we call pre-specified were specified in the script before that estimate was
computed, after related results had been seen; the model was written after most results, and no analysis was publicly
registered (\SIsec{sec:si-prereg}). Finally, the US analyses were re-run on a corrected 2026 release of the hiring
data, with a repaired name join, after all results had been seen (\SIsec{sec:si-data}).

Admission-cutoff designs with ranked application lists \citep{kirkeboen2016field, avery2013revealed} could ask
whether the return to a programme goes with its university's standing or with its department's
(\SIsec{sec:si-model}). Because a department's rank and its university's rank are highly correlated within fields
(\SItab{tab:si-map}), such a design would need many cutoffs to vary the department's rank at a given university rank.
The department-level events that public data let us study, first ABET or AACSB accreditations and Nobel prizes, are
followed by no large change in first-year pay or in completions by students already enrolled
(\SIsec{sec:si-events}).

% =====================================================================================================================
\section*{Methods}

\paragraph{Data.} All inputs are public; \SItab{tab:si-data} lists sources and releases, and \SItab{tab:si-flow} the
samples. The US academic ordering is the published SpringRank \citep{debacco2018physical} field rank $F$ and
university-wide rank $G$ of \citet{wapman2022quantifying}, estimated from a census of faculty at 368 PhD-granting US
institutions, 2011--2020, and signed so that higher means higher standing. We use the 2026 regeneration of their
data, the GitHub release at commit 841e9c9 (ref.~\citenum{wapman2026data}). An earlier public export \citep{wapman2022data} had
a truncated list of hiring edges, and an earlier version of this analysis used it (\SIsec{sec:si-data}). US pay is the median
earnings of bachelor's completers four years after completion in the College Scorecard field-of-study file
\citep{scorecard}, which covers completers who received federal (Title IV) aid and who work and are not enrolled.
Programmes are matched to departments by normalised institution name, with 65 hand-checked pairs for institutions whose
Scorecard names carry campus or city qualifiers, and by a curated map of 66 fields, 65 of which have a 2026 rank. The
canonical sample keeps the 55 fields with at least 15 matched universities, a threshold that keeps each field-level
correlation from resting on a handful of universities. The main gradient specification loses universities without a
reported SAT average, among them the University of California campuses, and this loss falls partly on highly ranked
departments (\SIsec{sec:si-slope}). Scorecard reports one set of medians per OPEID6 identifier, so campuses that share
one, such as Rutgers New Brunswick and Camden, carry the same pay (\SIsec{sec:si-data}). PSEO \citep{pseo, pseo2026}
gives median pay one, five and ten years after graduation for three-year graduation cohorts, which lets us follow the
same cohorts over their careers. We use release V4.14.1 and keep, within each field and cohort from 2001--2003 to
2010--2012, the institutions observed at all three horizons (40 fields, 127 institutions), so that a change across
horizons is not a change in the set of institutions. UK pay is the provider $\times$ subject median from the
Longitudinal Education Outcomes (LEO) file \citep{dfe_leo_provider}, one, three and five years after graduation, for
the cohorts 2013/14 to 2016/17. The file also gives medians by prior-attainment band, which let the UK analyses control
graduates' own prior attainment. The UK department measure is the REF 2021 grade-point average of the unit of
assessment matched to the subject \citep{ref2021results}, for 15 subjects; REF 2014 is used for the test--retest
reliability of the rating and for one variant of the gradient. Because the UK has no department-level hiring rank, the
UK $G$ is a SpringRank of 11,987 doctorate-to-faculty moves between institutions in Great Britain in ORCID-derived
mobility records \citep{li2026orcid, ror2026data}. It orders universities much as REF does (Spearman 0.815 across 131
universities; \SItab{tab:si-map}). UK intake is the share of a programme's graduates who entered with at least 360
UCAS tariff points; in the two-level contrast, suppressed cells of one or two graduates are coded as one. US
selectivity (mean SAT, admission rate) and Pell share come from the Scorecard institution files and describe the
university, not the programme.

\paragraph{Agreement at two levels.} To separate universities from departments, we compute one statistic at two
levels on the same programmes. Each ordering is converted to Blom normal scores within field (US) or subject (UK) and
standardised there, so that both orderings have the same distribution in every field and a linear correlation can be
pooled across fields. The across-university statistic is the Pearson correlation of these scores over all programmes.
It is a within-field total correlation, which mixes variation between universities with variation between the
departments of one university. The within-university statistic is the Pearson correlation of the scores' residuals on
university and field fixed effects. The university-mean statistic is the Pearson correlation, over universities, of
the university means of the field-demeaned scores (\SIsec{sec:si-contrast}). The same programmes enter all three:
3,961 US programmes (214 universities, 53 fields; the cells of the main gradient specification) and 864 UK programmes
(111 providers, 15 subjects). The US contrast that removes only fixed effects was added after the first results,
because the pre-specified within value also removed field-specific slopes on mean SAT, admission rate, Pell share and
$G$ that the across value did not; that stricter value is reported in \SIsec{sec:si-contrast}.

\paragraph{Reliabilities.} A test of the two levels needs each measure's reliability within universities, where the
noise is largest. For REF it is the within-provider correlation of REF 2014 and REF 2021 ratings, read as a
test--retest reliability, which real change between the exercises lowers. For pay it is one minus the share of
residual variance due to the sampling noise of each median, computed from the released quartiles and earner counts
(\SIsec{sec:si-data}). For intake it is an odd/even-cohort split half stepped up with the Spearman--Brown formula
\citep{spearman1910correlation, brown1910some}. For $F$, split halves of the public hiring edges give reliabilities
that average 0.89 across 58 fields. Because SpringRank on the complete public edges reproduces the published rank, this
is the sampling reliability of the published rank itself; an extrapolation to all hires with a US doctorate (0.92)
serves as a sensitivity. \SIsec{sec:si-contrast} describes the versions of the residual reliability of $F$ and which
analysis uses each.

\paragraph{Testing the two levels.} Our pre-specified rule divided each within correlation by the square root of the
product of the two reliabilities and compared the result with the across correlation. Simulation later showed that
this rule is not a calibrated test. Normal scores computed within subject lower the within correlation even without
noise, the intake split half overstates intake reliability, and in the US the rule corrects for noise in $F$ but not
in pay. Under equal agreement, the rule declares lower agreement within universities in up to 17\% of simulated UK
data sets for the normal scores used here and up to 53\% for linear scores, and in 81\% of US data sets, against a
nominal 5\%. We report its verdicts (\SIsec{sec:si-rule}) but do not use them as evidence. Instead, after the first
results had been seen, we calibrated the test by a parametric bootstrap. Latent scores are $T_k=u_k[i]+v_k[p]$ for
each ordering $k$, with a university component $u$, common to all fields of a university, and a programme component
$v$ on the real programmes and universities, $\mathrm{corr}(u)=\rho_{\mathrm{between}}$ and
$\mathrm{corr}(v)=\rho_{\mathrm{within}}$. Observed measures are generated as the data are made, so that the simulated
statistics face the same noise as the real ones: REF and $F$ carry classical normal error ($F$ at its field's
reliability), pay medians carry each programme's own sampling noise, and UK intake is a binomial count published with
LEO's rounding and suppression. Nuisance parameters are fitted on the same data by indirect inference to the observed
variance shares, the REF test--retest correlation and the intake moments; under the null, a common $\rho$ is
calibrated to the observed across correlation. Every simulated data set passes through the full pipeline, including
the fixed effects and the reliabilities. With 600 null data sets, the one-sided $P$ is
$(1+\#\{\text{null}\le\text{observed}\})/(1+600)$, so its smallest attainable value, $P=0.0017$, means that no
simulated within correlation is as low as the observed one. Model-based estimates free $\rho_{\mathrm{within}}$ and
match the observed within and across correlations; their approximate simulation intervals are in
\SItab{tab:si-model}.

\paragraph{Agreement in each field.} For field $f$, agreement is the Spearman correlation, across the universities
$i$ that teach $f$, between an academic ordering ($F_{if}$ or $G_i$) and median pay $Y_{if}$; the Supplementary Information
and the panels of Fig.~\ref{fig:career} call it coupling. A rank correlation compares the two orderings
without assuming a functional form, and it is not driven by the skew of pay medians. Partial agreement correlates the
rank residuals of the academic ordering and of $Y$ on the ranks of the controls. Means over fields weight fields
equally, so that large fields do not dominate the field map.

\paragraph{Department gradient.} The correlations measure how closely the orderings agree; the gradient measures how
much pay goes with a department's standing. The within-university slope is $\beta$ in $\log Y_{if}=\alpha_i+\mu_f+\beta
F^{z}_{if}+\sum_c\gamma_{fc}C_{ic}+\varepsilon_{if}$, with $F^{z}$ standardised within field and $C$ the
university-level controls of the main specification. The university fixed effects $\alpha_i$ restrict the comparison
to departments of the same university, and the field-specific slopes $\gamma_{fc}$ let selectivity and $G$ carry a
different weight in each field, as the field map indicates they do. Because the controls absorb most of $F$, the
reliability correction divides $\beta$ by the reliability of $F$ net of the design, with the public-edge and with the
extrapolated reliability of $F$. The UK slope uses REF rank in place of $F$, with provider and subject fixed effects.

\paragraph{Positive control and prediction.} The positive control correlates, within universities, the department
rank $F$ with research measures of the same departments from the National Research Council's assessment of
research-doctorate programmes (data for 2005--06) \citep{nrc2011data}. It uses the within statistic of the two-level
contrast on departments with a clean one-to-one match in fields that name the same discipline. The primary measure,
named in advance, is publications per faculty member; the NRC's overall rankings also weight doctoral production and
placement and serve only as secondary measures. The prediction analysis holds out whole universities in repeated
tenfold cross-validation and predicts each programme's log pay from field effects and university-level covariates
($G$, mean SAT, admission rate, Pell share, control and state earnings level), with and without $F$; a second task holds
out one programme at a time and adds the mean residual pay of the university's other programmes. No estimate uses the
held-out pay (\SIsec{sec:si-pcp}).

\paragraph{Career time.} Within each field and graduation cohort, PSEO cells keep the institutions observed at all
three horizons, with at least 15 per cell, so the same institutions are compared at one, five and ten years. The
career-time slope is the per-year change in mean agreement across horizons within a cohort. Such a change mixes
career time with calendar time, because a cohort is observed at later horizons in later years. Comparing cohorts at
matched calendar years bounds the career-time part, under the assumption that calendar and cohort effects are not
negative; a surface of graduation windows and horizons, on fields observed with the same institutions in every cell,
gives a second bound of the same kind (\SIsec{sec:si-career}). These bounds are identified sets under sign
restrictions, not confidence intervals. To ask which academic ordering gains weight, we regress pay ranks on the ranks
of $F$ and $G$ within field and cohort; the standardised coefficients are the loadings $b_F$ and $b_G$. Their per-year
changes are estimated cell by cell and, as a check, in a pooled form. The UK career slope uses LEO cohorts observed
one, three and five years after graduation at the same providers.

\paragraph{Inference.} Interval bars in the figures and brackets in the SI are 95\% intervals. Most are percentile
intervals from a cluster bootstrap that resamples universities (US) or providers (UK), because the programmes of one
university are not independent of one another. Two variants are used: means over UK subjects resample subjects, then
providers, and slopes compared across specifications reuse the same draws, so that changes are paired. The two-level
contrast uses 2,000 replicates and recomputes every reliability in each; the US slopes reuse the 999 replicates of
the original specification. PSEO career-time slopes, the audits of $F$ and the comparison of Scorecard with PSEO use
the two-way (field, university) cluster variance \citep{cameron2011robust}, because the same universities recur across
fields and cohorts. In Fig.~\ref{fig:gradient}, the REF 2014 slope has a cluster-robust interval, and the slope across
universities, whose source gives a cluster-robust SE, is drawn as a point. The grey bands of
Fig.~\ref{fig:levels} are ranges of simulated values, not intervals. The construct and model tests (\SIsec{sec:si-model}), the two-level rule, the
calibration audit, the positive control and the prediction rule were specified in the analysis scripts before the
corresponding estimate was computed, after related results had been seen, except the heterogeneous-loading nulls of the
calibration audit, chosen after a development run. They were not publicly registered. The calibrated simulation test, the
like-for-like US contrast and the production adjustment inside the gradient design were added after the first results
had been seen, and a few checks were computed after the main results (\SIsec{sec:si-prereg}).

\paragraph{Ethics.} The study is a secondary analysis of publicly available data: aggregate statistics released by
statistical agencies and government departments, and public records of researchers' academic careers (the
faculty-hiring data and ORCID-derived mobility records), which we analyse only in aggregate. No human participants
were recruited, no private or restricted records were accessed, and no information on identifiable individuals is
reported; ethical approval was therefore not sought.

\paragraph{Use of generative AI.} The author used Claude (Anthropic) to assist with collecting and processing the public data, writing and checking the analysis
code, and drafting and editing the manuscript. Each analysis was re-computed independently, in a separate session,
before its results were used. The author directed the work, reviewed the results and the text, and takes full
responsibility for the content.

\subsection*{Data availability}
All data analysed are public and are listed, with releases and access dates, in \SItab{tab:si-data} and in
\texttt{data/sources.csv} of the repository \url{https://github.com/soray42/universities-not-departments}: the
faculty-hiring ranks of \citet{wapman2026data}, College Scorecard field-of-study and institution files
\citep{scorecard}, PSEO earnings \citep{pseo2026}, the DfE LEO provider file \citep{dfe_leo_provider}, REF 2021 results
\citep{ref2021results} and REF 2014 results, ORCID-derived mobility records \citep{li2026orcid}, and the NRC data
table of research-doctorate programmes \citep{nrc2011data}, whose terms allow personal, noncommercial use only, so it is
not redistributed and no programme-level value from it is reported.

\subsection*{Code availability}
The analysis scripts, which are seeded and give
byte-identical outputs on re-run, and the figure scripts with their shared style module are in the public repository
\url{https://github.com/soray42/universities-not-departments}, which contains the code that reproduces every number,
table and figure from these inputs; raw files are downloaded from their sources and are not redistributed.
\SItab{tab:si-code} maps each section, table and figure to the script that produces it.

% =====================================================================================================================
% References: one list numbered in order of first citation, so Methods references continue the main-text numbering.
\bibliographystyle{unsrtnat}
\bibliography{refs_merged,refs_short,refs_nhb}

\section*{Acknowledgements}
I thank the creators of the faculty-hiring data for making them public and for releasing a corrected version, and the
U.S. Department of Education, the U.S. Census Bureau, the UK Department for Education, UK Research and Innovation
and the National Academies for publishing the data used here. This work received no specific funding.

\section*{Author contributions}
H.Y. conceived and designed the study, assembled and analysed the data, and wrote the manuscript.

\section*{Competing interests}
The author declares no competing interests.

\end{document}
